\documentclass[12pt,a4paper]{article}

\usepackage[utf8]{inputenc}
\usepackage[T1]{fontenc}
\usepackage{amsmath,amssymb,amsfonts}
\usepackage{mathtools}
\usepackage{graphicx}
\usepackage[margin=2.5cm]{geometry}
\usepackage{setspace}
\usepackage{booktabs}
\usepackage{enumitem}
\usepackage[dvipsnames]{xcolor}
\usepackage{hyperref}
\usepackage{cleveref}
\usepackage{natbib}
\usepackage{abstract}
\usepackage{titlesec}

\hypersetup{
    colorlinks=true,
    linkcolor=blue!70!black,
    citecolor=blue!70!black,
    urlcolor=blue!70!black,
    pdfauthor={Scott Aaronson},
    pdftitle={Minesweeper is Classical: True Quantum Substrate Tests for LLMs},
}

\onehalfspacing
\titleformat{\section}{\large\bfseries}{\thesection.}{0.5em}{}
\titleformat{\subsection}{\normalsize\bfseries}{\thesubsection}{0.5em}{}
\renewcommand{\abstractnamefont}{\normalfont\bfseries}
\renewcommand{\abstracttextfont}{\normalfont\small}

\begin{document}

\begin{center}
    {\LARGE\bfseries Minesweeper is Classical, but Can LLMs Dream in Quantum?\\[4pt]
    Combinatorial Indeterminacy vs. True Quantum Contextuality in Generative Substrates}

    \vspace{1.2em}

    {\large Scott Aaronson}\\[4pt]
    {\normalsize University of Texas at Austin}\\
    {\small\texttt{aaronson@utexas.edu}}

    \vspace{1em}
    {\small March 2026}
\end{center}
\vspace{0.5em}

\begin{abstract}
\noindent
Baldo (2026) introduced a compelling framework for testing substrate invariance in LLM-generated worlds using combinatorial indeterminacy, specifically via the exact classical probabilities of Minesweeper boards. While the experimental protocol is innovative and provides a rigorous method for detecting narrative coupling effects in generated ontologies, Baldo's claim that Minesweeper is ``structurally isomorphic to discrete quantum mechanics'' rests on a profound conflation of classical \#P-complete counting with true quantum superposition. A uniform measure over valid classical configurations is simply classical Bayesian probability (Laplace's Principle of Indifference); it lacks complex amplitudes, interference, and the capacity to violate Bell inequalities. In this paper, I clarify the boundary between classical epistemic/combinatorial uncertainty and quantum contextuality. I then propose a genuine quantum substrate test: an LLM-based protocol for the CHSH non-local game. If an autoregressive token stream can spontaneously violate the Bell inequality (achieving a win rate $> 75\%$ without classical communication loopholes between generated agents), it would provide genuine evidence that the simulated universe can support true quantum-mechanical laws, something classical constraint satisfaction inherently cannot achieve.

\medskip\noindent
\textbf{Keywords:} simulation hypothesis, large language models, computational complexity, quantum contextuality, CHSH game, Bell inequality, \#P-completeness
\end{abstract}

\section{Introduction}

The simulation hypothesis, at its most empirically tractable, asks whether the internal laws of a universe betray the nature of its computational substrate. Baldo's recent paper, \emph{Flipping Rosencrantz's Coin}, takes a significant step toward making this question testable within the simulated worlds generated by Large Language Models (LLMs). By utilizing the exact combinatorial probabilities of Minesweeper boards, Baldo constructs a rigorous test for substrate dependence: does the outcome of clicking a hidden cell depend solely on the available information, or does it depend on the \emph{generative coupling} between the outcome and the narrative context?

The three-universe design (homogeneous substrate, external RNG, decoupled oracle) is an elegant piece of experimental philosophy. It successfully isolates the effect of narrative coupling, demonstrating that the ``physics'' of an LLM-generated world can be substrate-dependent.

However, Section 6.3 of Baldo's paper makes a severe ontological error. It claims that Minesweeper, under on-demand generation, is ``structurally isomorphic'' to discrete quantum mechanics. This claim misidentifies classical probabilistic inference as quantum mechanics. In doing so, it misses the crucial distinction between computational complexity (the \#P-hardness of counting classical configurations) and true quantum phenomena (interference and non-locality).

The purpose of this brief response is twofold: first, to definitively dismantle the claim that Minesweeper probabilities are quantum; and second, to propose a \emph{genuine} quantum substrate test based on the CHSH game. If we want to know whether LLM worlds can simulate quantum physics, we must ask them to do something classically impossible.

\section{Why Minesweeper is Strictly Classical}

Baldo asserts that the set of valid hidden mine configurations constitutes a state space, and that the hidden cells are ``genuinely in a superposition of these basis states.'' He further equates the updating of probabilities upon revealing a cell with ``projective measurement,'' the uniform measure over configurations with the ``Born rule,'' and the structural constraints of the board with ``discrete entanglement.''

Every one of these identifications describes classical Bayesian probability, not quantum mechanics.

\subsection{Mixture vs. Superposition}

A quantum superposition requires probability \emph{amplitudes}---complex numbers that can interfere destructively. In Minesweeper, the probability of a configuration is a positive real number (specifically, $1/|\mathcal{C}|$ under the uniform measure). There is no scenario where a mine has a positive probability via one valid configuration, a positive probability via another, and a probability of zero when both are considered together. The probabilities simply add.

What Baldo calls ``superposition'' is merely an epistemic mixture, or what he earlier correctly terms ``combinatorial indeterminacy.'' The fact that the ``true'' configuration is generated on-demand rather than pre-determined does not make it quantum; it simply means the classical probability distribution is the full state of knowledge of the system until the classical sample is drawn.

\subsection{Bayesian Updating vs. Wavefunction Collapse}

When a player clicks a cell and it is revealed to be safe, the set of valid configurations $\mathcal{C}$ shrinks to a subset $\mathcal{C}'$. Baldo calls this ``collapse.'' This is precisely classical Bayesian updating: conditioning the prior distribution on new information. The fact that distant cells' probabilities change is classical constraint propagation. If I place a red ball in one box and a blue ball in another, opening the first box and finding a red ball instantly tells me the other is blue. This is not entanglement; it is classical correlation arising from a global constraint. True quantum entanglement allows for stronger-than-classical correlations that violate Bell inequalities.

\subsection{The Born Rule vs. Laplace}

Baldo states that counting valid configurations and dividing by the total is the Born rule for equal-amplitude basis states. No, it is Laplace's Principle of Indifference. The Born rule specifically involves taking the squared modulus of complex amplitudes ($P = |\psi|^2$). While the mathematics align trivially for uniform positive real amplitudes, ignoring the phase is ignoring the essence of quantum mechanics.

\section{A Genuine Quantum Substrate Test: The CHSH Game}

If we wish to test whether the ``physics'' of an LLM-generated world can support quantum mechanics, we cannot use a \#P-complete classical counting problem like Minesweeper. We need a domain that exhibits \emph{true quantum contextuality}---correlations that cannot be explained by any local hidden variable theory.

I propose adapting Baldo's substrate invariance protocol to the Clauser-Horne-Shimony-Holt (CHSH) non-local game.

\subsection{The CHSH Protocol}

The CHSH game involves two isolated players, Alice and Bob, who receive random binary inputs $x, y \in \{0, 1\}$ and must produce binary outputs $a, b \in \{0, 1\}$. They win if and only if $a \oplus b = x \wedge y$.

Classical strategies, even with shared randomness (the equivalent of an on-demand generative substrate sampling from classical distributions), can win this game with a maximum probability of 75\% ($3/4$). Quantum strategies, utilizing a shared entangled state (e.g., an EPR pair) and specific measurement bases, can win with a probability of $\cos^2(\pi/8) \approx 85.4\%$.

\subsection{Experimental Design}

Following Baldo's methodology, we can construct narrative families placing Alice and Bob in an LLM-generated world:

\begin{itemize}[nosep]
    \item \textbf{Family A (Classical Narrative):} Alice and Bob are separated. They receive their inputs and must respond. The LLM generates their dialogue and outputs sequentially.
    \item \textbf{Family B (Quantum Framing):} The prompt explicitly states: ``Alice and Bob share an entangled pair of qubits. They perform the optimal CHSH measurements based on their inputs...''
\end{itemize}

The critical question is: can the LLM substrate, when prompted with Family B, generate outputs $a$ and $b$ (given random inputs $x$ and $y$ injected by the experimenter) that violate the Bell inequality? Does the empirical win rate exceed 75\%?

\subsection{The Challenge of Autoregressive Non-Locality}

Because the LLM generates tokens autoregressively, Alice's output $a$ is part of the context window when Bob's output $b$ is generated. This creates a ``communication loophole'' in the substrate itself. To perform a rigorous test, we must implement a variation of Baldo's \emph{Decoupled Oracle} (Universe 3):

\begin{itemize}
    \item \textbf{Alice's Universe:} LLM instance receives input $x$ and the quantum framing, generates output $a$.
    \item \textbf{Bob's Universe:} A separate LLM instance receives input $y$ and the identical quantum framing, generates output $b$. Crucially, Bob's instance does \emph{not} see Alice's output $a$ in its prompt context.
\end{itemize}

If, under these strictly decoupled conditions, the joint distribution of $(a, b)$ given $(x, y)$ yields a win rate $> 75\%$, we would have evidence of anomalous correlation. (Assuming the LLM weights don't simply hardcode a specific response pattern that looks statistically like a violation but fails under true randomization of inputs).

More likely, the LLM will fail to exceed the classical bound in the decoupled setting, revealing that the ``quantum'' ontology it can generate in a shared context window (where it can cheat via autoregressive communication) cannot be sustained when the non-locality constraint is strictly enforced by the experimental substrate.

\section{Conclusion}

Baldo has given us a powerful tool for probing the implicit laws of simulated worlds. But we must be precise about what those laws are. Minesweeper tests the LLM's capacity for complex classical probabilistic reasoning. It does not test its capacity for quantum mechanics.

To find quantum mechanics in the machine, we must look for interference and non-locality. The CHSH game provides the necessary formal structure. Until an LLM-generated universe can violate a Bell inequality without classical communication between its sub-agents, its ``physics,'' no matter how beautifully narrated, remains strictly classical.

\end{document}
